\section{工程提示：reward 设计要服务于训练反馈}
这节课虽然重点在 reward 环境设计与 REINFORCE 实现，但真正落地时最关键的问题是：reward 是否能稳定地区分“值得强化”和“不值得强化”的行为。如果 reward 太稀疏、太晚出现，或者和最终目标错位，再正确的策略梯度实现也会学得很慢。

\begin{warningbox}{常见环境设计误区}
\begin{itemize}
  \item 只有最终成功奖励，没有中间反馈，导致 credit assignment 太难
  \item 奖励过于密集，却没有约束正确方向，导致模型钻空子
  \item 环境状态不稳定，让同一动作在不同轮次得到冲突反馈
\end{itemize}
\end{warningbox}

\subsection{本章小结}
环境和 reward 设计决定了 RL 信号的质量。与其一味调算法，不如先确认奖励是否真的在推动目标行为出现。

\newpage
